\section{Fourth Principle: Revise the Loop That Revises the Work}

A decision changes the environment. New evidence changes the standard. A changed standard alters the next decision. Improvement therefore includes the procedure used to judge improvement.

\subsection{A Recursive Process Can Learn or Drift}

For a scalar iteration,
\begin{equation}
 x_{n+1}=g(x_n),
\end{equation}
a continuously differentiable map with a fixed point $x^\ast$ is locally attracting when $|g'(x^\ast)|<1$. In several dimensions, the Jacobian's eigenvalues govern the corresponding local result under the usual regularity assumptions.

An institution updating its evaluation rules is more complicated. Objectives, participants, evidence, and external conditions can all change. The equation supplies a reason to inspect the update rule, rather than a promise of convergence.

The useful question is whether a revision preserves learning. Can errors be identified? Can affected people challenge the interpretation? Is there enough continuity to distinguish a better rule from a lucky outcome?

\subsection{Keep Some Practice Inside the Person}

Assistance can make a task cheaper while moving practice elsewhere. A person who wants to retain a capability needs some occasions to exercise it and obtain interpretable feedback.

Choose the capability first. Someone preserving spatial judgment might sometimes navigate without turn-by-turn instructions. Someone learning a language might draft a passage before requesting corrections. Someone responsible for code should retain the ability to read, test, and explain the parts they will maintain.

The useful unit is the learning loop, rather than unaided effort for its own sake. A good tutor, simulator, or assistant can improve that loop. Deliberately withholding an effective aid is worthwhile only when it serves a defined purpose.

Defaults can protect practice without requiring a fresh contest of will at every moment. A regular session, a shared commitment, or a tool configuration can make a desired activity ordinary. Waiting can preserve a decision point; feedback can be delayed when immediate scoring distorts the activity.

Different people have different resources. Time, disability, care responsibilities, money, and fatigue change which arrangements are feasible. Assistance often expands agency most for people facing the highest barriers. Preserving a skill should not become a demand that everyone accept the same avoidable burden.

Organizations influence the available choices. A worker cannot preserve every valuable learning process in unpaid hours while being evaluated only on immediate output. Training, staffing, and assessment have to carry part of the cost.

\subsection{Build Organizations That Produce Warrants}

A laboratory, hospital, engineering firm, court, auditor, or certification body can produce a judgment that others are entitled to rely on. Its product is partly a warrant: a result connected to evidence, procedure, expertise, and responsibility.

The name alone supplies none of this. An approval chain can generate signatures without generating confidence. An audit can become a ritual. A university can reward output volume while neglecting understanding.

Three conditions make checking substantive. Reviewers need access to relevant original evidence. Their incentives must allow an unfavorable finding. Their authority must be strong enough to change the proposed action.

Tools can accelerate each part. They can organize evidence, check consistency, reproduce calculations, and highlight anomalies. A trustworthy organization captures those gains while examining the new failure modes introduced by the tools.

Four slow assets deserve deliberate support: informative exposure over time, a record of trustworthy behavior, a clear distribution of duties, and the independence needed to challenge a profitable error.

The timescales vary. Some checks can be automated almost completely. Some clinical or material observations require time to pass. An organization should distinguish a queue caused by poor administration from a follow-up period that supplies necessary information.

Markets can reward dependable judgment, but information asymmetry, delayed consequences, and shared benefits can leave safeguards underprovided. Public standards, professional norms, procurement, liability, and independent institutions can change those incentives.

An organization that invests in reliability may also save money by preventing rework and incidents. The cost is not necessarily a permanent handicap. The comparison must include the losses the faster-looking arrangement produces later.

Records connect the periods. A finding should remain traceable after staff turnover or a model update. A failed prediction should enter the next review rather than disappear behind a new version.

\subsection{A Life Has More Than One Axis}

Output is measurable and often useful. It is an incomplete account of a person's life.

National accounts measure defined economic aggregates. They do not claim to measure every form of worth. Trouble begins when a convenient output measure becomes the sole basis for recognizing activity.

AARP's March 2026 report estimates that unpaid adult caregiving in the United States supplied 49.5 billion hours in 2024, with an imputed economic value of about \$1.01 trillion \cite{aarp2026valuing}. The valuation uses a method and wage assumptions; it is not money actually paid. The report also changed aspects of its method, so simple comparisons with older headline totals can mislead.

The estimate makes visible work already occurring. It also exposes the difference between being socially necessary and having a funded position.

Function and experience are separate evaluative dimensions. A highly productive activity can leave a person depleted. An activity with modest market output can sustain a valuable relationship or a meaningful experience. Calling the dimensions separate does not establish statistical independence between them.

The same person can value both. Automation can free time, improve access, and remove drudgery. The benefit becomes part of a good life through what the person can then do, experience, and sustain.

This is why productivity need not be defended by finding ever more output to assign to everyone. A society can use increased capacity to support care, education, leisure, participation, and a more secure floor beneath choice.

The fourth principle is to keep revising the arrangement by its actual consequences. Preserve learning, use tools to improve evidence, fund the necessary slow work, and evaluate lives on more than one axis.
